\documentclass[11pt]{article}

% Packages
\usepackage{amsmath,amssymb,amsthm} % Math stuff
\usepackage{enumitem} % For improved lists
\usepackage{fancyhdr} % For headers
\usepackage{geometry} % Page geometry
\usepackage{graphicx}
\usepackage{booktabs}
\usepackage{subcaption}
\geometry{margin=1in}
\usepackage{xcolor}
\usepackage{float}
\usepackage{mathtools}
\usepackage[colorlinks=true, 
    linkcolor=blue, 
    citecolor=blue, 
    urlcolor=blue, 
    filecolor=blue]{hyperref}

% Header
\pagestyle{fancy}
\fancyhf{}
\lhead{Asher Baraban, Matej Cerman \\ Martin Hiti}
\chead{14.383}
\rhead{Problem Set 1}
\cfoot{\thepage}

\newtheorem{theorem}{Theorem}[section]

% For easier problem/solution environments
\newenvironment{problem}[1]
    {\bigskip\noindent\textbf{Problem #1.}\par\setlength{\parindent}{0pt}}
    {\medskip}
\newenvironment{solution}[1]
    {\bigskip\noindent\textbf{Solution #1.} \par}
    {\medskip}

\newcommand\barbelow[1]{\stackunder[1.2pt]{$#1$}{\rule{.8ex}{.075ex}}}

% Citations
\usepackage{natbib}
\bibliographystyle{unsrtnat}

% Math shortcuts
\newcommand{\ev}{\mathbb{E}}
\newcommand{\evn}{\mathbb{E}_n}
\newcommand{\R}{\mathbb{R}}

% Title
\title{Problem Set 1 \\ 14.383}
\author{Asher Baraban, Matej Cerman, Martin Hiti}
\date{\today}

% Document
\begin{document}


\maketitle

\noindent \textit{Notes}:
\begin{itemize}
    \item All code is in \url{https://github.com/ABaraban/14.383_pset4}
    \item All three group members took 14.381 last semester
\end{itemize}

\subsection*{Chapter 9} 
\begin{problem}{9.1 (Asher)}
Implement standard panel data estimators (fixed effects, first
differences, and GMM-FD2 approaches) for the first or third
empirical example. These are implemented in standard software
such as Stata or R. Very briefly explain the assumptions
you need to make for these estimators to be consistent for
causal effects, and why these assumptions may or may not
hold in these examples. Very briefly explain how you are
taking care of clustering.
\end{problem}

\begin{solution}{}
Below, we implemented the Democracy example. For our estimates to be causal we need
\[\epsilon_{it} \perp (X_i^t, a_i), X_i^t = (D_{it}, ..., D_{i1}, Y_{i(t-1)}, ..., Y_{i1})\]
In words, democracy and past GDP are orthogonal to current and future GDP shocks. We also need the error terms to be serially uncorrelated conditional on sufficient lags of GDP. We are clustering by country 

\begin{table}[htbp]
\centering
\begin{tabular}{rrrrrr}
  \hline
 & Pooled & FD & FE & GMM-FD1 & GMM-FD2 \\ 
  \hline
Democracy & 0.46 & 0.94 & 1.89 & 4.45 & 3.91 \\ 
  CSE & 0.26 & 1.20 & 0.65 & 2.77 & 1.70 \\ 
  BSE & 0.24 & 1.25 & 0.67 & 2.93 & 1.57 \\ 
  L1.log(gdp) & 1.33 & 0.33 & 1.15 & 1.30 & 1.00 \\ 
  CSE1 & 0.05 & 0.05 & 0.05 & 0.14 & 0.07 \\ 
  BSE1 & 0.05 & 0.05 & 0.05 & 0.14 & 0.07 \\ 
  L2.log(gdp) & -0.18 & 0.17 & -0.12 & -0.15 & -0.08 \\ 
  CSE2 & 0.06 & 0.04 & 0.06 & 0.28 & 0.06 \\ 
  BSE2 & 0.06 & 0.04 & 0.06 & 0.22 & 0.06 \\ 
  L3.log(gdp) & -0.11 & 0.06 & -0.07 & -0.43 & -0.04 \\ 
  CSE3 & 0.05 & 0.02 & 0.04 & 0.28 & 0.04 \\ 
  BSE3 & 0.05 & 0.03 & 0.04 & 0.25 & 0.04 \\ 
  L4.log(gdp) & -0.05 & -0.05 & -0.08 & 0.18 & -0.08 \\ 
  CSE4 & 0.03 & 0.04 & 0.02 & 0.14 & 0.03 \\ 
  BSE4 & 0.03 & 0.04 & 0.02 & 0.14 & 0.03 \\ 
  LR-Democracy & 288.03 & 1.91 & 16.05 & 46.75 & 20.91 \\ 
  CSE5 & 162.89 & 2.45 & 6.67 & 52.61 & 10.60 \\ 
  BSE5 & 157.69 & 2.54 & 6.44 & 36.11 & 9.16 \\ 
  J-test &  &  &  & 70.71 & 130.23 \\ 
  p-val &  &  &  & 0.00 & 1.00 \\ 
  dof &  &  &  & 31.00 & 463.00 \\ 
   \hline
\end{tabular}
\end{table}
\end{solution}

\newpage

\subsection*{Chapter 10}


\begin{problem}{10.1 (Martin)}
Explain how the bias can arise in high dimensional problems using a simple example.
\end{problem}

\begin{solution}{10.1}
Under standard low/fixed dimensional GMM, we have shown that the difference between out estimated parameter and the true value of the parameter is of order $1/\sqrt{n},$ which vanishes as $n \to \infty.$ However, when the problem becomes high dimensional---either because the dimensionality of the parameter is high or the number of moments is high---we are not always guaranteed such consistency.

To see this, we present a contrived example that shows how bias can arrive in high dimensional problems. Partition the parameter vector in to a target parameter $\alpha$ and a high dimensional nuisance parameter $\gamma$:
\begin{align*}
    \theta = (\alpha', \gamma')' \qquad \theta_0 = (\alpha'_0, \gamma'_0)'
\end{align*}
where subscripts. We approximate the case with a high dimensional nuisance parameter as $p = \dim(\gamma) \to \infty $ as $n\to \infty$ but a fixed dimension for the target parmeter $d_\alpha = \dim(\alpha).$

Now, suppose that we are interested in the target parameter is $\alpha = \gamma'\gamma$ which is one dimensional ($d_{\alpha} = 1)$ and that the true value satisfies $\alpha_0 = \gamma_0'\gamma_0 = 1.$ Suppose that we have an estimator $\hat{\gamma}$ of $\gamma_0$ that satisfies
\begin{align*}
    G_n \coloneq \hat{\gamma} - \gamma_0 \sim \mathcal{N}(0,I_p/n)
\end{align*}
That is, the estimation error in each component of $\hat{\gamma}$ is mean zero, independent, and of order $1/\sqrt{n}$ (so it goes to 0 as $n \to \infty$). Consider the plug-in estimator of $\alpha_0$: $\hat{\alpha} = \hat{\gamma}'\hat{\gamma}.$ This estimator obeys an exact Taylor expansion
\begin{align*}
    \hat{\alpha} - \alpha_0 
    &=\hat{\gamma}'\hat{\gamma} - \gamma_0'\gamma_0\\
    &= (\gamma_0 + G_n )'(\gamma_0 +G_n) - \gamma_0'\gamma_0 \\
    &= \gamma_0'\gamma_0 + 2\gamma_0'G_n + G_n'G_n - \gamma_0'\gamma_0 \\
    &=  2\gamma_0'G_n + G_n'G_n \\
    &= \frac{2\sqrt{n}\gamma_0'G_n}{\sqrt{n}} + G_n'G_n \tag{$\star$}
\end{align*}
Let $Z_n \coloneq 2\gamma_0'G_n\sqrt{n}.$ Note that  $\gamma_0'G_n$ is a linear combination of the components of $G_n.$ Since $G_n \sim \mathcal{N}(0,I_p/n),$ any linear combination of $G_n$ is also normal. In particular, 
\begin{align*}
    \gamma_0'G_n \sim \mathcal{N}(0,\gamma_0'\frac{I_p}{n}\gamma_0) = \mathcal{N}(0,1)
\end{align*}
since $\gamma_0'\gamma_0 = 1.$ Therefore 
\begin{align*}
    \frac{Z_n}{\sqrt{n}} = \frac{2\sqrt{n}\gamma_0'G_n}{\sqrt{n}} = 2\gamma_0'G_n \sim 2\mathcal{N}(0,1) \sim \mathcal{N}(0,4)
\end{align*}
Thus, the first term in the expansion $(\star)$ is mean zero and asymptotitcally normal, so it is not the source of any bias. Now, let's consider the quadratic term $G'_nG_n$ Let's decompose this term into a deterministic part (bias) and a remainder part:
\begin{align*}
    G_n'G_n = \ev[G_n'G_n ] + (G_n'G_n - \ev[G_n'G_n ])
\end{align*}
Since $G_n \sim \mathcal{N}(0,I_p/n)$ we have that $\frac{1}{\sqrt{n}}G_n\sim\mathcal{N}(0,I_p).$ Therefore $nG_nG_n = \frac{1}{\sqrt{n}}G_n\frac{1}{\sqrt{n}}G_n$ is the sum of $p$ independent squared standard normals $\implies nG_nG_n \sim \chi^2_p.$ Thus, 
\begin{align*}
    \ev[nG_n'G_n ] = p \implies \ev[G_n'G_n ] = \frac{p}{n}
\end{align*}
Thus, 
\begin{align*}
    G_n'G_n = \frac{p}{n} + \underbrace{\left(G_n'G_n - \frac{p}{n}\right)}_{=:r_n}
\end{align*}
Since $nG_nG_n \sim \chi^2_p$, the central limit theorem as $p \to \infty$ allows us to find the asymptotic distribution of $r_n \coloneq G_n'G_n - \frac{p}{n}:$
\begin{align*}
    &\frac{nG_n'G_n - p}{\sqrt{2p}} \overset{a}{\sim} \mathcal{N}(0,1) \\
    \implies&\frac{nG_n'G_n - p}{\sqrt{p}} \overset{a}{\sim} \mathcal{N}(0,2) \\
    \implies&r_n = G_n'G_n - \frac{p}{n} \sim\mathcal{N}(0,2)\sqrt{p}/n
\end{align*}
Thus, $r_n = O_p(\frac{\sqrt{p}}{n}).$

Putting it all together, we can write the full expansion as:
\begin{align*}
    \hat{\alpha} - \alpha_0 = \underbrace{\frac{Z_n}{\sqrt{n}}}_{\mathcal{N}(0,4/n)} + \underbrace{\frac{p}{n}}_{\text{bias}} + \underbrace{r_n}_{O_p(\sqrt{p}/n)}
\end{align*}
where $Z_n \sim \mathcal{N}(0,4)$, the bias term is $b/n = p/n$, and $r_n = O_p(\sqrt{p}/n)$. To assess when the normal approximation $\sqrt{n}(\hat{\alpha} - \alpha_0) \overset{a}{\sim} \mathcal{N}(0,4)$ is valid, rescale by $\sqrt{n}$:
\begin{align*}
    \sqrt{n}(\hat{\alpha} - \alpha_0) = \underbrace{Z_n}_{\mathcal{N}(0,4)} + \underbrace{\frac{p}{\sqrt{n}}}_{\sqrt{n} \cdot \text{bias}} + \underbrace{O_p\!\left(\frac{\sqrt{p}}{\sqrt{n}}\right)}_{\sqrt{n}\cdot r_n}
\end{align*}
For the normal approximation to hold, both the bias term and the remainder must vanish asymptotically:
\begin{align*}
    \frac{p}{\sqrt{n}} \to 0 \iff \frac{p^2}{n} \to 0 \qquad \text{(binding constraint)}\\
    \frac{\sqrt{p}}{\sqrt{n}} \to 0 \iff \frac{p}{n} \to 0 \qquad \text{(weaker condition)}
\end{align*}
The bias term is the binding constraint. The source of the bias is the nonlinearity of the plug-in estimator $\hat{\alpha} = \hat{\gamma}'\hat{\gamma}$: even though each component of $\hat{\gamma}$ is unbiased for the corresponding component of $\gamma_0$, squaring and summing across all $p$ dimensions accumulates the estimation variance into a systematic upward shift of size $p/n$. Each of the $p$ dimensions contributes a bias of $1/n$, so the total bias grows with the dimension. When $p$ is small relative to $n$ this is negligible, but as $p$ grows the accumulated bias becomes non-negligible relative to the $1/\sqrt{n}$ standard deviation of the estimator, invalidating the standard normal approximation even if $\hat{\alpha}$ remains consistent.
\end{solution}



\begin{problem}{10.2}
Replicate (approximately) the empirical results of Table 10.1
using standard software such as R or Stata. [Note that you
will not be able to replicate exactly the split-sample bias corrections
and bootstrap standard errors because they depend
on the software and random seed used.]
\end{problem}

\begin{solution}{10.2}
We present results using the provided code in the table below. The sample contains $n=147$ countries and $T=18$ years for both GMM1 and GMM2. The GMM1 model is the Anderson-Hsiao estimator, which uses fewer moment conditions ($m=200$), while GMM2 is the Arellano-Bond estimator that takes advantage of all restrictions implied by sequential exogeneity, so $m=632$. There are $p=169$ parameters, with parameters of interest being coefficients on democracy and the lags of GDP (the dependent variable), and incidental parameters being country and time effects. As a result, $m^2/n$ (which is the likely bottleneck) is large for in GMM1 and very large in GMM2, and we need to worry about bias due to dimensionality.

For each model, we show three columns: the raw outputs without bias correction, the bias-corrected estimates (BC) which use a single sample split, and BC5 which averages over 5 random sample splits.

The CSE rows report analytical standard errors of the GMM estimators clustered at the country level, while BSE rows report standard errors from a bootstrap with 500 repetitions.

Finally, the last section computes the long-run effects on GDP of a transition to democracy, not estimated directly but computed as a transformation of the other coefficients.

\begin{center}
    % latex table generated in R 4.5.1 by xtable 1.8-8 package
% Thu Apr 16 11:54:56 2026
\begin{table}[ht]
\centering
\begin{tabular}{rcccccc}
  \hline
 & GMM1 & GMM1-BC & GMM1-BC5 & GMM2 & GMM2-BC & GMM2-BC5 \\ 
  \hline
Democracy & 4.45 & 4.21 & 5.61 & 3.91 & 4.57 & 5.13 \\ 
  CSE & 2.77 &  &  & 1.70 &  &  \\ 
  BSE & 3.00 & 3.87 & 4.07 & 1.72 & 2.10 & 2.15 \\ 
  L1.log(gdp) & 1.30 & 1.49 & 1.47 & 1.00 & 1.05 & 1.04 \\ 
  CSE1 & 0.14 &  &  & 0.07 &  &  \\ 
  BSE1 & 0.17 & 0.25 & 0.19 & 0.07 & 0.13 & 0.11 \\ 
  L2.log(gdp) & -0.15 & -0.36 & -0.27 & -0.08 & -0.10 & -0.08 \\ 
  CSE2 & 0.28 &  &  & 0.06 &  &  \\ 
  BSE2 & 0.18 & 0.28 & 0.24 & 0.06 & 0.08 & 0.08 \\ 
  L3.log(gdp) & -0.43 & -0.58 & -0.54 & -0.04 & -0.04 & -0.04 \\ 
  CSE3 & 0.28 &  &  & 0.04 &  &  \\ 
  BSE3 & 0.21 & 0.34 & 0.29 & 0.03 & 0.07 & 0.05 \\ 
  L4.log(gdp) & 0.18 & 0.39 & 0.30 & -0.08 & -0.08 & -0.09 \\ 
  CSE4 & 0.14 &  &  & 0.03 &  &  \\ 
  BSE4 & 0.15 & 0.21 & 0.20 & 0.03 & 0.05 & 0.04 \\ 
  LR-Democracy & 46.75 & -3.12 & 44.31 & 20.91 & 24.41 & 25.34 \\ 
  CSE5 & 52.61 &  &  & 10.60 &  &  \\ 
  BSE5 & 36.29 & 59.96 & 51.62 & 8.64 & 10.90 & 11.42 \\ 
   \hline
\end{tabular}
\end{table}

\end{center}

The results show that the bias corrections meaningfully change the estimates, particularly for GMM2, which likely suffers from a larger bias problem. At the same time, the bootstrap standard errors largely go up for the bias-corrected estimates, and are by-and-large larger than analytical standard errors, which suggests that analytical SEs might not be capturing the additional errors introduced by estimation in a small $T$ (relative to $n$) panel.

The estimates of long-run effects obtained from GMM1 are larger, but are also much less precisely estimated, highlighting the increase in statistical power from using more over-identifying restrictions in GMM2. The bias-corrected GMM2 estimate suggests that a permanent transition to democracy raises GDP by $~20\%$

\end{solution}



\end{document}